% =====================================================================
% chapters/01_Introduction.tex — LaTeX rendering of ../sections/01_introduction.en.md (2026-09-24 00:01:41)
% Section 1 of the SHORT paper. \input from main.tex, first of the body.
% This file DEFINES \label{sec:intro}. It carries no \ref: the former 1.4 Organisation,
%   with its six renvois and \label{sec:organisation} (referenced nowhere), was removed from
%   the master on 2026-09-23 and is removed here on 2026-09-24.
% Citations: lammer2026gta (GTA), vu2025modeling (the architecture this work follows),
%   tisseo2023emc2 (the survey; the acronym EMC^2 is deliberately not written in the body,
%   the citation carries it).
% Derived file: regenerate after every validated pass of ../sections/01_introduction.en.md.
%   The "source:" comments of the Markdown stay there and are not carried over.
% =====================================================================

\section{Introduction}
\label{sec:intro}

\subsection{What generative agents promise}
\label{sec:promise}

Existing multi-agent simulations of multimodal urban mobility rely on tabular models
estimated from household travel surveys, or on explicit rules defined by domain experts. A
household travel survey asks the residents of one territory to describe every trip they made
on one day. In both, the space of representable behaviours is fixed \emph{a priori}: a factor
encoded as neither variable nor rule cannot influence any decision.

The fixed behaviour space binds precisely where simulation is asked to help. A city tests a
transport policy on a simulated population before it builds anything. The answer it gets is
bounded by the behaviour its model can represent.

Generative agents built on language models promise to lift that bound. Fed on
massive corpora, they carry decision heuristics that surveys do not record, and they adapt
without explicit rules. A strike, a heatwave or a news item about safety in the metro could
then reach a decision that no survey column holds. They also promise to work without the
local calibration data that agent-based models require. Detailed multimodal trajectory
datasets exist for a few large metropolitan areas only.

\subsection{Measuring the promise against a real population}
\label{sec:gap}

Published evaluations of these agents rarely have a human population to compare with. Most
judge a generative agent against a rule-based agent, or against aggregate patterns the
authors produced themselves.

GTA \citep{lammer2026gta} confronts a simulated modal split with a national travel
survey. The modal split is the share of trips made by each mode. GTA draws its personas from
that survey. A language model writes each of their day plans. Lacking a baseline, it
compares itself with the observed splits of other German states. Copying Hamburg's shares scores a
root-mean-square error of 1.99 on Berlin, against 4.07 for its own simulation.

This paper supplies four things missing from these comparisons. A baseline beneath the
agents says what a decision-maker reaches with no behavioural knowledge of the territory.
Models fitted on that territory's own survey say what its data support. The same
inputs on both sides make a gap attributable to the decision-maker rather than to what it was
told. A reading below the aggregate says whether two decision-makers that produce the same
shares decide alike.

\subsection{What this paper does}
\label{sec:contributions}

This paper measures where verbalised deliberation earns its place in a mobility agent.
Verbalised deliberation is the model writing its reasoning in words before it answers. In a
population of mobility agents, it does not improve fidelity to the ordinary regime the
survey describes. It brings in events that no variable encodes. We trace the path of one such
event to the decision.

We measure on one study area, the 453 communes around Toulouse that its Cerema-certified 2023
household travel survey covers \citep{tisseo2023emc2}. A controlled cohort of 1{,}000 synthetic personas makes the
3{,}299 trips of the day we score. Fifteen decision-makers run on those trips. We call
decision-maker anything that turns a trip description into a probability over the options
offered. The agents live in a GAMA multi-agent simulation of the study area, with its real
networks and timetables. That simulation follows the GAMA--OpenTripPlanner--LLM architecture
of \citet{vu2025modeling}.

The paper rests on three contributions. The first is a comparison bench with equal inputs
between generative agents and tabular models, on a synthetic cohort controlled against that
survey. The second is the position of fifteen decision-makers on that bench, and two
dissociations the aggregate shares do not show. The third is the traced path of a
non-tabulated event to the decision, inside a generative agent with memory.

